\documentclass[10pt,a4paper,ssfamily]{exam}
\usepackage[brazil]{babel}
\usepackage[numbers,sort&compress]{natbib}
\usepackage[utf8]{inputenc}
\renewcommand{\baselinestretch}{1.0}
\usepackage{epsfig}
\usepackage{mhchem}
\usepackage{graphicx}
\usepackage{makeidx}
\usepackage{multicol}
\usepackage{multirow}
\usepackage{amssymb}
\usepackage{amsmath}
\DeclareMathOperator{\sen}{sen}
\newcommand{\listfont}{\bf\fontencoding{OT1}\sffamily\large}
\newcommand{\titlesfont}{\fontencoding{OT1}\sffamily\large}
\renewcommand{\baselinestretch}{1.0}
\newcommand{\Mg}{Mg$^{2+}$}
\renewcommand{\t}{\textrm}
\newcommand{\cc}[1]{[\textrm{#1}]}
\usepackage{enumitem}
\setlist{nosep, topsep=2pt, partopsep=0pt, parsep=0pt, itemsep=2pt}
\newcommand{\bmat}{\left[\begin{array}{cc}}
\newcommand{\emat}{\end{array}\right]}
\newcommand{\nomera}{
  \noindent
  \begin{tabular}{|p{0.7\textwidth}|p{0.3\textwidth}|}
  Nome: & RA: \\
  \hline
  \end{tabular}
}

\newcommand{\qini}{\begin{enumerate}[resume]\item}
\newcommand{\qend}{\end{enumerate}}

\newcommand{\oC}{$^\circ$C}
\newcommand{\kcalmol}{~kcal~mol$^{-1}$}
\newcommand{\moll}{~mol~L$^{-1}$}
\newcommand{\molkg}{~mol~kg$^{-1}$}
\newcommand{\gl}{~g~L$^{-1}$}
\newcommand{\e}[1]{$\times 10^{#1}$}

\newcommand{\eps}{\varepsilon}

\newcommand{\Resposta}[1]{\vspace{0.5cm}\noindent{\bf Resposta:}\\ #1
\begin{center}\rule{\textwidth}{0.4pt}\end{center}}
\renewcommand{\Resposta}[1]{}
\newcommand{\resposta}[1]{\vspace{0.5cm}\noindent{\bf Resposta:}\\ #1
\begin{center}\rule{\textwidth}{0.4pt}\end{center}}


\begin{document}
\sffamily
\pagestyle{empty}
\nomera

\begin{center}
{\bf QG101 - Química Geral}\\
{\bf Aula 7 - Ligações Covalentes e Estruturas de Lewis}\\
\underline{Justifique suas respostas.}
\end{center}

\vspace{0.2cm}
\begin{center}{\bf Questões}\end{center}

% ---------------------------------------------------------------
\qini
\textbf{Estruturas de Lewis e carga formal.}
Considere as seguintes espécies: \ce{H2O}, \ce{NH3}, \ce{CO2} e \ce{N2}.

\begin{enumerate}
  \item Para cada espécie, conte o total de elétrons de valência e
        desenhe a estrutura de Lewis completa, indicando todos os pares
        ligantes e não-ligantes.

  \item Usando a fórmula
        \[
          \text{CF} = (\text{e}^-\ \text{valência})
                    - (\text{e}^-\ \text{não-ligantes})
                    - \tfrac{1}{2}(\text{e}^-\ \text{ligantes})
        \]
        calcule a carga formal de cada átomo na estrutura de Lewis do
        \ce{CO2}. Verifique que a soma das cargas formais é igual
        à carga total da molécula.

  \item Uma estrutura alternativa para o \ce{CO2} seria
        \ce{:O^{-}-C#O+:}
        (uma ligação simples O--C e uma ligação tripla C--O).
        Calcule as cargas formais nessa estrutura e explique por que
        ela é \textbf{menos estável} que \ce{O=C=O}.
\end{enumerate}

\vspace{0.3cm}
\qend

% ---------------------------------------------------------------
\qini
\textbf{Ressonância no íon carbonato (\ce{CO3^{2-}}).}

\begin{enumerate}
  \item Conte o total de elétrons de valência do \ce{CO3^{2-}}
        (lembre-se de incluir os dois elétrons da carga negativa).

  \item Desenhe as \textbf{três} estruturas de ressonância do
        \ce{CO3^{2-}}, indicando as cargas formais de cada átomo
        em cada estrutura.

  \item Calcule a \textbf{ordem de ligação média} C--O usando:
        \[
          \text{Ordem média} =
          \frac{\text{soma das ordens nas estruturas de ressonância}}
               {\text{número de estruturas equivalentes}}
        \]

  \item O comprimento experimental da ligação C--O no \ce{CO3^{2-}}
        é 129 pm. Sabendo que uma ligação simples \ce{C-O} tem
        143 pm e uma dupla \ce{C=O} tem 123 pm, verifique se esse
        valor é consistente com a ordem de ligação calculada.
        \textit{Sugestão}: estime o comprimento esperado por
        interpolação linear entre os dois valores de referência.
\end{enumerate}

\vspace{0.3cm}
\qend

% ---------------------------------------------------------------
\qini
\textbf{Geometria molecular e polaridade (VSEPR).}
Para cada espécie abaixo, aplique a teoria VSEPR para determinar
a geometria molecular e indicar se a molécula é polar ou apolar.

\begin{center}
\begin{tabular}{llcc}
\hline
Espécie & Fórmula & Pares ligantes & Pares livres (central)\\
\hline
Metano       & \ce{CH4}  & \phantom{0} & \phantom{0} \\
Amônia       & \ce{NH3}  & \phantom{0} & \phantom{0} \\
Água         & \ce{H2O}  & \phantom{0} & \phantom{0} \\
Dióx. de carbono & \ce{CO2} & \phantom{0} & \phantom{0} \\
Trifluoreto de boro & \ce{BF3} & \phantom{0} & \phantom{0} \\
\hline
\end{tabular}
\end{center}

\begin{enumerate}
  \item Preencha a tabela com o número de pares ligantes e
        não-ligantes no átomo central de cada espécie e determine
        a geometria molecular e o ângulo de ligação esperado.

  \item Os ângulos de ligação observados são:
        \ce{CH4}: 109,5°; \ce{NH3}: 107,0°; \ce{H2O}: 104,5°.
        Todos possuem 4 pares de elétrons ao redor do átomo central,
        mas os ângulos diminuem na sequência acima.
        Explique essa tendência usando o princípio de repulsão entre
        pares eletrônicos.

  \item Para quais das espécies listadas o vetor resultante dos
        momentos de dipolo de ligação é \textbf{nulo}? Justifique
        com base na geometria.
\end{enumerate}

\vspace{0.3cm}
\qend

\end{document}
